If $y=u=1$, then $x+z=2a$ and $xz=a-5$. The discriminant $4(a^2-a+5)$ is positive for every $a$, giving two ordered solutions.

If $y=u=-1$, the first and third equations force $a=0$. Then $x+z=0$, $xz=-1$, giving two more solutions only at $a=0$.

Otherwise $u=\dfrac{1}{y}$ and $y\ne\pm1$. Solving the first and third equations gives
\[
x=\dfrac{2ay}{1+y},\qquad z=\dfrac{2a}{1+y}.
\]
With $v=y+\dfrac{1}{y}$, the second equation becomes
\[
v^2+(5-a)v+4a^2-2a+6=0.
\]
Its discriminant is $D=(1-5a)(1+3a)$. Thus real $v$ require $-\dfrac{1}{3}\le a\le\dfrac{1}{5}$. For these parameters, both candidate values $\dfrac{(a-5\pm\sqrt D)}{2}$ are at most $-2$, since $D\le(1-a)^2$, with equality only at $a=0$. For $a\ne0$ they are strictly below $-2$, each giving two reciprocal real values of $y$.

Inside the interval, except at $0$, this gives four additional solutions; at its endpoints, two. At $a=0$, the $v$ values are $-3,-2$: the first gives two additional solutions, while $v=-2$ was already handled as $y=u=-1$. Altogether the answer is
\[
\boxed{\begin{cases}
6,&-\dfrac{1}{3}<a<\dfrac{1}{5},\\
4,&a=-\dfrac{1}{3}\text{ or }a=\dfrac{1}{5},\\
2,&a<-\dfrac{1}{3}\text{ or }a>\dfrac{1}{5}.
\end{cases}}
\]
